\chapter*{Введение}
\label{sec:intro}
\addcontentsline{toc}{chapter}{Введение}

Современный подбор персонала характеризуется высокой нагрузкой на этапе первичного отбора соискателей: в корпоративной практике на одну вакансию могут приходиться сотни откликов соискателей на вакансии на платформах/системах по поиску работы, что делает инструменты автоматизированного ранжирования и фильтрации соискателей ключевым элементом инфраструктуры процесса подбора персонала \cite{src02}.

Классическим программным контуром управления подбором персонала выступают системы автоматизированного подбора персонала (Applicant Tracking System), которые обеспечивают хранение цифровых профилей соискателей, то есть структурированных цифровых карточек соискателей с резюме, контактными данными, опытом работы, профессиональными навыками, статусом/этапом процесса их рассмотрения на вакансию и историей взаимодействия соискателя с вакансией в системе, управление стадиями воронки подбора персонала начиная с отклика соискателя на вакансию и заканчивая его оформлением в компании, где была опубликована вакансия, а также выдачу упорядоченных списков кандидатов для рассмотрения сотрудником по подбору персонала \cite{src02}.

На практике задача ранжирования кандидатов в системах автоматизированного подбора персонала формулируется как задача информационного поиска и рекомендаций: в качестве запроса выступает описание вакансии или формализованные требования к кандидату, его должностные обязанности, а документами являются резюме соискателей или их цифровые профили, которые необходимо отсортировать по степени соответствия требованиям вакансии. Под релевантностью в этой постановке понимается прикладная мера того, насколько опыт, навыки, уровень позиции и профессиональный контекст соискателя согласуются с требованиями конкретной вакансии и позволяют рекомендовать профиль к первичному рассмотрению специалистом по подбору персонала \cite{src02, src03}.

Кроме того, современные исследования подчеркивают важность различения справедливой экспозиции (равномерного распределения видимости кандидатов на различных позициях списка) и справедливых исходов (группы кандидатов, реально отобранных в шорт-лист). Корректировка алгоритма ранжирования сама по себе может не гарантировать справедливого результата, поскольку медиатором выступают пользовательские интерфейсы систем автоматизированного подбора персонала, стратегии просмотра списков соискателей и когнитивный эффект позиционного смещения (position bias), при котором пользователь непропорционально часто просматривает и выбирает элементы верхних позиций списка независимо от их абсолютного качества. Внедрение специализированных механизмов справедливого ранжирования результатов (fair ranking) является предметом будущих исследований и не рассматривается в данной работе \cite{src02}.

Цель работы заключается в проектировании модуля семантического сопоставления текстов вакансии и резюме на основе семантических представлений: модуль должен преобразовывать текстовое описание вакансии и резюме в числовые векторные представления, сравнивать их по смысловой близости, учитывать особенности предметной области, например за счёт тонкой настройки на ручной разметке, полученной от сотрудников по подбору персонала, и формировать упорядоченный список соискателей для первичного рассмотрения рекрутером.

Задачи работы:

\begin{compactenum}
  \item Разработать архитектуру интеграции модуля семантического сопоставления текстов вакансии и резюме в систему автоматизированного подбора персонала.
  \item Определить воспроизводимую методику оценки качества ранжирования списков соискателей по нормализованной дисконтированной кумулятивной выгоде (Normalized Discounted Cumulative Gain, NDCG), точности на первых K позициях (Precision at K, Precision@K), полноте на первых K позициях (Recall at K, Recall@K) и средней обратной позиции первого релевантного результата (Mean Reciprocal Rank, MRR).
  \item Рассчитать временной и стоимостной эффект предварительного ранжирования списков кандидатов моделью раздельного кодирования (bi-encoder) перед более дорогим экспертным оцениванием, например сотрудником по подбору персонала, или оцениванием большой языковой моделью.
\end{compactenum}

Методологическая база исследования включает современные нейросетевые модели контекстных представлений, в том числе архитектуру двунаправленных кодировочных представлений (Bidirectional Encoder Representations from Transformers), модель векторных представлений предложений на базе сиамских кодировочных сетей (Sentence Embeddings using Siamese Bidirectional Encoder Representations from Transformers Networks), подходы на базе архитектуры раздельного кодирования текстов (bi-encoder), а также трехступенчатые индустриальные схемы поиска соискателей с отбором, включающие предварительное ранжирование и окончательное переранжирование \cite{src01, src06, src08, src17}.

Актуальность работы состоит в том, что системы автоматизированного подбора персонала обрабатывают большие потоки резюме и нуждаются в воспроизводимых механизмах предварительного ранжирования кандидатов, которые уменьшают ручную нагрузку специалиста по подбору персонала и при этом оставляют кадровое решение под контролем человека.

Объектом исследования являются процессы автоматизированного сопоставления вакансий и цифровых профилей соискателей в системах подбора персонала.

Предметом исследования являются методы семантического представления текстов вакансий и резюме, алгоритмы вычисления их смысловой близости и программная интеграция такого ранжирования в контур системы автоматизированного подбора персонала.

Научная новизна работы заключается в адаптации архитектуры раздельного кодирования вакансии и резюме к доменному корпусу исторических кадровых событий и в сопоставлении качества такого подхода с классическими и нейросетевыми стратегиями ранжирования на фиксированной тестовой выборке.

Теоретическая значимость работы состоит в формализации задачи сопоставления текстов вакансии и резюме как задачи семантического ранжирования с использованием плотных векторных представлений, метрик информационного поиска и ограничений высокорисковой предметной области подбора персонала.

Практическая значимость работы заключается в разработке воспроизводимого программного прототипа, который может быть подключен к тестовому контуру системы автоматизированного подбора персонала Reqcore и использоваться для предварительной сортировки кандидатов перед экспертной проверкой специалистом по подбору персонала.

Содержание работы по разделам организовано следующим образом. В первой главе анализируются предметная область подбора персонала, методы обработки естественного языка и подходы к обучению ранжированию. Во второй главе рассматриваются исходные данные, их профилирование, очистка и разметка, требования к обработке текстов и фактическая интеграционная поверхность прототипа. В третьей главе проектируются и реализуются программный модуль, конвейер машинного обучения, сервисный программный интерфейс и подключение к Reqcore. В четвертой главе приводятся экспериментальная оценка, проверка интеграции, экономическое обоснование и валидация стратегий ранжирования.
